% manuscrito_final.tex -- AgroMoreira (SGA)
% Revista objetivo: Requirements Engineering (Springer Nature), plantilla sn-jnl.cls.
% Enfoque 2, legal-first (Amaral et al., 2021).
% Resultados, Discusion y Conclusiones redactados a partir de las salidas reales de
% 07_Datos/scripts/analisis_legalfirst.py sobre 07_Datos/datos_crudos/cobertura_legal.csv
% (26 criterios legales, McNemar chi2=16.056 corregido por continuidad, p<0.001, b=18, c=0).

\documentclass[pdflatex,sn-mathphys-num]{sn-jnl}

\usepackage{graphicx}
\usepackage{booktabs}
\usepackage{amsmath}
\usepackage{array}
\usepackage[T1]{fontenc}
\usepackage[utf8]{inputenc}

\begin{document}

\title[Cobertura de requisitos legales mediante un enfoque legal-first]{Requisitos legales de protecci\'on de datos y trazabilidad agroindustrial: un enfoque legal-first en Ecuador}

\author*[1]{\fnm{Jeanpierre Robinson} \sur{Espinoza}}\email{jrobinsone@uteq.edu.ec, ORCID 0009-0005-3302-1822}
\author[1]{\fnm{Danela Dayana} \sur{Arteaga \'Alava}}\email{darteagaa@uteq.edu.ec, ORCID 0009-0006-0318-1531}
\author[1]{\fnm{Kamila Annabella} \sur{Calle Delgado}}\email{kcalled@uteq.edu.ec, ORCID 0009-0002-8249-9601}
\author[1]{\fnm{Mar\'ia del Rosario} \sur{Escudero Plaza}}\email{mescuderop@uteq.edu.ec, ORCID 0009-0007-3212-9924}
\author[1]{\fnm{Roselyn Andreina} \sur{S\'anchez Centeno}}\email{rsanchezc4@uteq.edu.ec, ORCID 0009-0008-7204-8448}
\author[1]{\fnm{Gleiston Cicer\'on} \sur{Guerrero Ulloa}}\email{gguerrero@uteq.edu.ec, ORCID 0000-0001-5990-2357}

\affil*[1]{\orgdiv{Facultad de Ciencias de la Computaci\'on}, \orgname{Universidad T\'ecnica Estatal de Quevedo}, \city{Quevedo}, \country{Ecuador}}

\abstract{\textbf{Contexto:} Los requisitos legales de protecci\'on de datos y trazabilidad agroindustrial suelen quedar insuficientemente cubiertos cuando la elicitaci\'on se apoya solo en m\'etodos convencionales, sobre todo en pa\'ises en desarrollo con marcos legales recientes. \textbf{Objetivo:} Este estudio investiga qu\'e requisitos legales de la Ley Org\'anica de Protecci\'on de Datos Personales del Ecuador y de la normativa de trazabilidad agroindustrial quedan sin cubrir por la elicitaci\'on convencional, y c\'omo el enfoque legal-first de Amaral et al. extiende esa cobertura. \textbf{M\'etodo:} Mediante un estudio de caso en una finca ecuatoriana de policultivo, se construy\'o un modelo conceptual de cumplimiento legal con 26 criterios en tres bloques: protecci\'on de datos, trazabilidad del cacao y bioseguridad fitosanitaria. Se evalu\'o su cobertura por 39 requisitos funcionales y 21 no funcionales elicitados mediante 17 entrevistas, y se aplic\'o el enfoque legal-first para derivar ocho requisitos adicionales directamente del texto legal. \textbf{Resultados:} La cobertura convencional cubri\'o 7 de los 26 criterios (26,9\,\%) y subi\'o a 25 de 26 (96,2\,\%) tras aplicar el enfoque legal-first, una diferencia de 69,2 puntos porcentuales (IC 95\,\% bootstrap [50,0\,\%, 84,6\,\%]) estad\'isticamente significativa (McNemar con correcci\'on de continuidad, $\chi^2=16,06$, $p<0,001$), sin p\'erdida de cobertura en ning\'un criterio. \textbf{Conclusiones:} El enfoque legal-first extiende su validaci\'on m\'as all\'a del RGPD europeo a un marco legal latinoamericano reciente y a un dominio agroindustrial de policultivo subrepresentado en la literatura; el vac\'io que cierra se concentra en el marco legal m\'as reciente y menos incorporado a la pr\'actica cotidiana del sector.}

\keywords{Legal requirements, Traceability, Cacao industry, Data protection law, Requirements coverage}

\maketitle

\section{Introducci\'on}\label{sec:introduccion}

Agr\'icola Moreira, una finca ecuatoriana dedicada al cultivo simult\'aneo de cacao (\textit{Theobroma cacao}) y pl\'atano verde (\textit{Musa} spp.), gestiona hoy sus parcelas, labores, cosechas e insumos mediante bit\'acoras f\'isicas y hojas de c\'alculo, lo que dificulta el control de producci\'on, la detecci\'on temprana de enfermedades como la moniliasis y la sigatoka, y el cumplimiento de obligaciones legales recientes, la Ley Org\'anica de Protecci\'on de Datos Personales del Ecuador \cite{lopdp2021}, su Reglamento General \cite{rlopdp2023} y los requisitos de trazabilidad agroindustrial de AGROCALIDAD \cite{agrocalidad183}. El proyecto especifica los requisitos de un Sistema de Gesti\'on Agr\'icola con Inteligencia Artificial que centraliza ese registro e incorpora alertas con confirmaci\'on humana obligatoria. La elicitaci\'on convencional tiende a capturar necesidades funcionales expl\'icitas del negocio, pero puede dejar sin cubrir obligaciones legales que ning\'un informante menciona de forma espont\'anea, una brecha cr\'itica en un dominio agroindustrial de un pa\'is en desarrollo, donde el marco de protecci\'on de datos es reciente y la trazabilidad regulatoria apenas se est\'a formalizando.

Amaral et al. \cite{amaral2021} proponen un modelo conceptual para verificar el cumplimiento de requisitos de procesamiento de datos frente al RGPD europeo, y muestran que un enfoque legal-first, derivar requisitos de forma sistem\'atica a partir del texto legal, extiende la cobertura respecto de la elicitaci\'on convencional. Ese trabajo se valida sobre el RGPD y en contextos europeos, sin extenderse a marcos legales latinoamericanos ni a dominios agroindustriales. Torre et al. \cite{torre2021} y Negri-Ribalta et al. \cite{negriribalta2024} confirman que la investigaci\'on de cumplimiento en Ingenier\'ia de Requerimientos se concentra en el RGPD y carece de evidencia primaria de casos reales fuera de Europa. La trazabilidad del cacao, por su parte, se ha abordado como problema t\'ecnico de registro \cite{arsyad2018,marchese2022}, sin an\'alisis de la cobertura de requisitos legales. Ning\'un estudio localizado mide de forma emp\'irica la cobertura de requisitos legales de protecci\'on de datos y de trazabilidad agroindustrial en un caso latinoamericano, ni compara esa cobertura antes y despu\'es de aplicar un m\'etodo legal-first.

Este trabajo aporta un modelo conceptual verificable de cumplimiento legal para un sistema agroindustrial ecuatoriano, con 26 criterios en tres bloques normativos, una medici\'on emp\'irica de la cobertura de esos criterios antes y despu\'es de aplicar el enfoque legal-first de Amaral et al., y evidencia de un caso agroindustrial latinoamericano de policultivo, subrepresentado en la literatura de Ingenier\'ia de Requerimientos sobre cumplimiento legal. La lista definitiva de contribuciones se enuncia sobre la evidencia de la secci\'on de resultados.

La pregunta de investigaci\'on que gu\'ia el estudio es la siguiente. \textit{\'Que requisitos legales de la Ley Org\'anica de Protecci\'on de Datos y de trazabilidad agroindustrial no quedan cubiertos por los requisitos funcionales elicitados con m\'etodos convencionales, y c\'omo el enfoque legal-first de Amaral et al. extiende esa cobertura.}

\medskip
\noindent\textbf{Pregunta de investigaci\'on en formato PICOC.}
\begin{itemize}
\item \textbf{Population (poblaci\'on):} sistemas de gesti\'on agr\'icola con requisitos de cumplimiento legal en fincas agroindustriales latinoamericanas de policultivo, bajo un marco legal de protecci\'on de datos reciente.
\item \textbf{Intervention (intervenci\'on):} el enfoque legal-first de Amaral et al. \cite{amaral2021}, que deriva requisitos de forma sistem\'atica a partir de un modelo conceptual del texto legal.
\item \textbf{Comparison (comparaci\'on):} la cobertura de requisitos legales lograda mediante elicitaci\'on convencional, entrevistas semiestructuradas y cuestionario, sin aplicar el m\'etodo legal-first.
\item \textbf{Outcome (resultado):} la proporci\'on de los 26 criterios legales verificables de protecci\'on de datos y trazabilidad agroindustrial cubiertos por los requisitos funcionales, antes y despu\'es de aplicar el m\'etodo.
\item \textbf{Context (contexto):} un estudio de caso \'unico en Agr\'icola Moreira, una finca ecuatoriana de policultivo de cacao y pl\'atano verde, bajo la Ley Org\'anica de Protecci\'on de Datos Personales del Ecuador y la normativa de AGROCALIDAD.
\end{itemize}

\section{Trabajo relacionado}\label{sec:trabajo-relacionado}

La revisi\'on sigue el esquema abreviado de Kitchenham y Charters \cite{kitchenham2007}. Se ejecutaron tres cadenas de b\'usqueda en Scopus, IEEE Xplore y Google Scholar. La primera combina los t\'erminos de cumplimiento legal, \texttt{legal}, \texttt{regulatory}, \texttt{compliance}, \texttt{GDPR}, \texttt{data protection} y \texttt{privacy}, con los de Ingenier\'ia de Requerimientos. La segunda combina \texttt{traceability} y \texttt{supply chain} con \texttt{agri-food}, \texttt{agriculture}, \texttt{cocoa} y \texttt{cacao} y con \texttt{information system}, \texttt{requirements} y \texttt{blockchain}. La tercera combina \texttt{large language model}, \texttt{LLM} y \texttt{generative AI} con Ingenier\'ia de Requerimientos y con \texttt{quality}, \texttt{evaluation} y \texttt{coverage}. Se incluyeron estudios revisados por pares, con datos emp\'iricos o con un artefacto de modelo, m\'etodo o prototipo evaluado, publicados en conferencias o revistas de Ingenier\'ia de Requerimientos, sistemas de informaci\'on o trazabilidad agroalimentaria. Se excluyeron editoriales, res\'umenes de una p\'agina y textos sin evaluaci\'on ni datos. De la b\'usqueda se retuvieron quince estudios para el cuadro comparativo por su cercan\'ia directa a la pregunta de investigaci\'on. El n\'umero de registros recuperados por cada cadena y el n\'umero de estudios revisados a texto completo se reportan aqu\'i con la fecha de corte de la b\'usqueda sistem\'atica.

El Cuadro~\ref{tab:relacionado} organiza los quince estudios en tres bloques. El primero re\'une el enfoque legal-first y el cumplimiento normativo en Ingenier\'ia de Requerimientos, que es la base metodol\'ogica directa. El segundo re\'une la trazabilidad agroalimentaria y de cacao, que es el dominio del caso. El tercero re\'une la IA generativa y la calidad en la elicitaci\'on, que da contexto a la evoluci\'on del enfoque descrita en la Secci\'on~\ref{sec:metodologia}.

\begin{table}[htbp]
\caption{Cuadro comparativo del trabajo relacionado}\label{tab:relacionado}
\begin{tabular}{@{}p{2.2cm}p{1.9cm}p{2.6cm}p{2.6cm}p{3.4cm}@{}}
\toprule
Referencia & Tipo de estudio & Poblaci\'on e intervenci\'on & Resultados principales & Diferencia con nuestro trabajo \\
\midrule
\multicolumn{5}{@{}l}{\textit{Bloque A. Enfoque legal-first y cumplimiento normativo}} \\
\addlinespace
Breaux y Ant\'on \cite{breaux2008}, 2008 & M\'etodo con estudio de caso & Texto de la HIPAA Privacy Rule; extracci\'on manual y sistem\'atica de derechos y obligaciones de acceso a partir del articulado & Cobertura a nivel de enunciado de todo el documento regulatorio, con reglas para resolver referencias cruzadas y ambig\"uedad & Marco legal estadounidense de salud, no la protecci\'on de datos ecuatoriana ni la trazabilidad agroindustrial \\
\addlinespace
Massey et al. \cite{massey2010}, 2010 & M\'etodo con estudio de caso, sistema iTrust & Requisitos de seguridad y privacidad de un sistema de historia cl\'inica frente a HIPAA y HITECH; evaluaci\'on de conformidad legal de una especificaci\'on ya redactada & Cataloga los retos pr\'acticos de especificar para cumplimiento y clasifica los requisitos seg\'un su alineaci\'on con la ley & Eval\'ua una especificaci\'on existente, no compara la cobertura antes y despu\'es de un m\'etodo legal-first, y no toca el dominio agr\'icola \\
\addlinespace
Deng et al. \cite{deng2011}, 2011 & Marco metodol\'ogico con caso ilustrativo & Requisitos de privacidad de sistemas intensivos en software; aplicaci\'on de LINDDUN, un an\'alisis de amenazas de privacidad por categor\'ias & Deriva requisitos de privacidad y orienta la selecci\'on de tecnolog\'ias de protecci\'on a partir de las amenazas identificadas & Trabaja sobre amenazas de privacidad por dise\~no, no sobre la cobertura de obligaciones de un marco legal nacional concreto \\
\addlinespace
Amaral et al. \cite{amaral2021}, 2021 & Modelo conceptual con evaluaci\'on & Obligaciones de procesamiento de datos del RGPD; conceptualizaci\'on basada en modelos para verificar el cumplimiento de una especificaci\'on & Modelo de cumplimiento con reglas verificables, base del enfoque legal-first que adopta este estudio & Validado sobre el RGPD y en contexto europeo, sin extensi\'on a un marco legal latinoamericano ni a la trazabilidad agroindustrial \\
\addlinespace
Torre et al. \cite{torre2021}, 2021 & Modelo conceptual con aplicaci\'on y reporte de experiencia & Articulado del RGPD; construcci\'on de un modelo UML con restricciones OCL y su aplicaci\'on a casos reales & Modelo del RGPD con trazabilidad al articulado, 35 reglas de cumplimiento y 20 puntos de variaci\'on & Modelo gen\'erico del RGPD, no instanciado para la LOPDP del Ecuador ni combinado con normativa agr\'icola \\
\addlinespace
Negri-Ribalta et al. \cite{negriribalta2024}, 2024 & Mapeo sistem\'atico de 90 art\'iculos entre 2016 y 2022 & Literatura que aborda el RGPD desde la Ingenier\'ia de Requerimientos; clasificaci\'on y s\'intesis de m\'etodos, t\'ecnicas y vac\'ios & Muestra que la mayor\'ia de propuestas son parciales y que falta evidencia emp\'irica de casos reales & Estudio secundario centrado en el RGPD, sin evidencia primaria de un caso agroindustrial ni de otro marco legal nacional \\
\addlinespace
\multicolumn{5}{@{}l}{\textit{Bloque B. Trazabilidad agroalimentaria y de cacao}} \\
\addlinespace
Arsyad et al. \cite{arsyad2018}, 2018 & Dise\~no de artefacto con prototipo & Cadena de suministro de cacao y su documentaci\'on de respaldo; blockchain de dos factores que enlaza la documentaci\'on legal con el registro de trazabilidad mediante marca de agua digital & Un mecanismo que permite comprobar la integridad conjunta de la documentaci\'on y del rastro de trazabilidad & Se centra en la integridad criptogr\'afica del registro, no en la elicitaci\'on ni en la cobertura de requisitos legales \\
\addlinespace
Marchese y Tomarchio \cite{marchese2022}, 2022 & Dise\~no de artefacto con prototipo & Cadena de suministro agroalimentaria gen\'erica; sistema de trazabilidad sobre blockchain con contratos inteligentes para registrar los eventos de la cadena & Un prototipo que gestiona y consulta los eventos de trazabilidad de forma descentralizada & Propuesta t\'ecnica de trazabilidad, sin an\'alisis de requisitos legales ni trabajo de campo con productores \\
\addlinespace
\multicolumn{5}{@{}l}{\textit{Bloque C. IA generativa y calidad en la elicitaci\'on}} \\
\addlinespace
Krishna et al. \cite{krishna2024}, 2024 & Estudio comparativo & Ingenieros de nivel inicial y modelos GPT-4 y CodeLlama; redacci\'on de una especificaci\'on de requisitos completa por cada fuente & Calidad comparable a la de un ingeniero junior y mucho menos tiempo con el modelo & Sin evaluaci\'on a ciegas por expertos independientes, caso universitario y no agroindustrial \\
\addlinespace
Cheng et al. \cite{cheng2026}, 2026 & Revisi\'on sistem\'atica & Literatura sobre IA generativa en Ingenier\'ia de Requerimientos; s\'intesis del estado del campo & La fase de gesti\'on de requisitos recibe atenci\'on insuficiente en la literatura & Estudio secundario, no evidencia emp\'irica primaria ni de dominio agr\'icola \\
\addlinespace
Arora et al. \cite{arora2024}, 2024 & Art\'iculo de posici\'on & Pr\'actica profesional de Ingenier\'ia de Requerimientos; discusi\'on estructurada del rol de los modelos de lenguaje & Un mapa de oportunidades y riesgos del uso de modelos de lenguaje en RE & No mide de forma emp\'irica la calidad ni la cobertura de los requisitos \\
\addlinespace
Vogelsang y Fischbach \cite{vogelsang2024}, 2024 & Gu\'ia sistem\'atica & Tareas de procesamiento de lenguaje natural dentro de RE; formulaci\'on de directrices de uso & Un procedimiento recomendado para aplicar modelos de lenguaje a esas tareas & De alcance general, sin foco en la cobertura de requisitos legales \\
\addlinespace
Vogelsang y Borg \cite{vogelsang2019}, 2019 & Estudio cualitativo con 4 entrevistas & Cient\'ificos de datos que construyen sistemas de aprendizaje autom\'atico; entrevistas semiestructuradas sobre sus necesidades de requisitos & La explicabilidad aparece como un nuevo requisito de calidad en sistemas de IA & Perspectiva de quienes construyen el sistema, no del cumplimiento normativo \\
\addlinespace
Ferrari et al. \cite{ferrari2016}, 2016 & Estudio experimental & Analistas de requisitos en entrevistas de elicitaci\'on; an\'alisis de transcripciones & La ambig\"uedad y el conocimiento t\'acito persisten aun con analistas expertos & Explica por qu\'e la elicitaci\'on convencional deja vac\'ios, sin proponer un m\'etodo legal-first \\
\addlinespace
Chazette y Schneider \cite{chazette2020}, 2020 & Estudio conceptual & Explicabilidad como requisito no funcional; definici\'on del concepto y cat\'alogo de recomendaciones & Un marco para tratar la explicabilidad como requisito no funcional & Se enfoca en la explicabilidad, no en la cobertura de requisitos legales \\
\botrule
\end{tabular}
\end{table}

El enfoque legal-first est\'a validado sobre el RGPD europeo y en contextos de salud o software gen\'erico \cite{breaux2008,massey2010,amaral2021,torre2021,negriribalta2024}, y la trazabilidad de cacao se ha abordado como problema t\'ecnico de registro \cite{arsyad2018,marchese2022}. Ning\'un estudio localizado mide la cobertura de requisitos legales de protecci\'on de datos y de trazabilidad agroindustrial en un caso agroindustrial latinoamericano de policultivo de cacao y pl\'atano verde, ni compara esa cobertura antes y despu\'es de aplicar un enfoque legal-first sobre un marco legal ecuatoriano. Este estudio ocupa ese espacio.

\section{Metodolog\'ia}\label{sec:metodologia}

Esta secci\'on resume el protocolo experimental registrado del estudio.

\subsection{Dise\~no del estudio}

El estudio combina un dise\~no cualitativo, con entrevistas y observaci\'on, y una comparaci\'on de proporciones pareadas sobre 26 unidades. Se construy\'o un modelo conceptual de cumplimiento legal con 26 criterios de cumplimiento, C1 a C26, en tres bloques normativos, la Ley Org\'anica de Protecci\'on de Datos Personales del Ecuador, la Resoluci\'on 183 de AGROCALIDAD sobre buenas pr\'acticas y trazabilidad del cacao, y la Resoluci\'on 0072 sobre bioseguridad y manejo fitosanitario. Cada criterio queda enlazado a su art\'iculo legal verificado.

Para cada criterio se eval\'ua la cobertura, como variable binaria de cubierto o no cubierto, en dos momentos. Primero por el conjunto de requisitos funcionales y no funcionales elicitado mediante m\'etodos convencionales. Despu\'es tras aplicar el m\'etodo legal-first de Amaral et al. \cite{amaral2021}, que deriva requisitos adicionales directamente del texto legal, en este caso ocho requisitos, RF-22, RF-23, RF-24, RF-25, RF-36, RF-37, RF-38 y RF-39 de la especificaci\'on de requisitos del sistema.

\subsection{Participantes y reclutamiento}

La recolecci\'on de datos de campo combina 17 entrevistas semiestructuradas con personas distintas, de las cuales un subconjunto sirvi\'o adem\'as como sesi\'on de validaci\'on con walkthrough del prototipo (correspondencia EV$\leftrightarrow$ENTR documentada en el ERS/SRS), y un cuestionario aplicado a 114 personas, con 68 respuestas del perfil dominante (agricultor), por encima del m\'inimo de 60 previsto en el protocolo. El levantamiento del cuestionario y su reconciliaci\'on con el m\'inimo planificado se documentan en \texttt{07\_Datos/desviaciones.md}. El consentimiento informado se obtuvo conforme a la Ley Org\'anica de Protecci\'on de Datos Personales del Ecuador \cite{lopdp2021}. La regla de decisi\'on para marcar un criterio como cubierto se document\'o por escrito antes de evaluar, y la asignaci\'on se contrasta con al menos un segundo revisor independiente.

\subsection{Modelo de cumplimiento legal y m\'etodo legal-first}

El modelo conceptual sigue la l\'ogica de Amaral et al. \cite{amaral2021}. A partir del articulado de cada norma se identifican obligaciones verificables y se expresan como criterios de cumplimiento con un enunciado medible. Un criterio se marca como cubierto solo si existe al menos un requisito funcional o no funcional, con criterio de aceptaci\'on medible, que lo satisface. Los criterios que quedan sin ning\'un requisito que los cubra son los vac\'ios legales que responden a la pregunta de investigaci\'on, y los requisitos derivados del texto legal para cerrarlos son la extensi\'on de cobertura que aporta el m\'etodo legal-first.

\subsection{Plan de an\'alisis}

El plan de an\'alisis se registra en OSF \cite{nosek2018} antes de evaluar formalmente la cobertura. Incluye la tabla de cobertura pareada de cubierto o no cubierto para los 26 criterios en los dos momentos, estad\'isticos descriptivos de cobertura globales y desglosados por bloque normativo y por fuente de evidencia, la comparaci\'on de proporciones pareadas mediante prueba de McNemar, por tratarse de datos binarios pareados sobre las mismas 26 unidades, y el intervalo de confianza del 95 por ciento de la diferencia de proporciones por bootstrap con 10\,000 r\'eplicas. El nivel de significancia es 0{,}05. Se documentan dos desviaciones respecto del ideal metodol\'ogico. La primera, que parte del trabajo de campo de elicitaci\'on comenz\'o antes del registro formal en OSF, como trabajo exploratorio de entregas anteriores del curso. La segunda, que el enfoque metodol\'ogico pas\'o por dos formulaciones previas antes de fijarse en el dise\~no legal-first, por asignaci\'on expl\'icita de la gu\'ia oficial de la entrega. Ninguna decisi\'on de an\'alisis se tom\'o a partir de resultados preliminares.

\section{Resultados}\label{sec:resultados}

El an\'alisis de cobertura se ejecuta con \texttt{07\_Datos/scripts/analisis\_legalfirst.py} sobre los 26 criterios legales de \texttt{07\_Datos/datos\_crudos/cobertura\_legal.csv}, con la salida verificada en \texttt{07\_Datos/resultados/tabla\_mcnemar.csv} y \texttt{07\_Datos/resultados/descriptivos\_bloque.csv}, evaluados en dos momentos: antes de aplicar el m\'etodo legal-first, con los requisitos elicitados mediante las 17 entrevistas, incluido el subconjunto con sesi\'on de walkthrough del prototipo, y despu\'es, tras derivar directamente del texto legal los ocho requisitos adicionales RF-22, RF-23, RF-24, RF-25, RF-36, RF-37, RF-38 y RF-39.

\subsection{Cobertura global}

La cobertura convencional cubre 7 de los 26 criterios (26,9\,\%). Tras aplicar el m\'etodo legal-first, la cobertura sube a 25 de 26 criterios (96,2\,\%), una diferencia de 69,2 puntos porcentuales (intervalo de confianza del 95\,\% por bootstrap con 10\,000 r\'eplicas, [50,0\,\%, 84,6\,\%]). La comparaci\'on pareada de proporciones mediante la prueba de McNemar, con correcci\'on de continuidad, resulta estad\'isticamente significativa, $\chi^2(1)=16,06$, $p<0,001$.

La Tabla~\ref{tab:mcnemar} presenta la tabla de contingencia $2\times2$ de los pares de cobertura. De los 26 criterios, 18 pasaron de no cubiertos a cubiertos al aplicar el m\'etodo legal-first, 7 estaban cubiertos en ambos momentos y ninguno retrocedi\'o, de cubierto a no cubierto. S\'olo un criterio permanece sin cobertura en ambos momentos.

\begin{table}[htbp]
\caption{Tabla de contingencia pareada de cobertura, antes y despu\'es del enfoque legal-first ($n=26$ criterios)}\label{tab:mcnemar}
\begin{tabular}{@{}lcc@{}}
\toprule
 & Legal-first: no cubierto & Legal-first: cubierto \\
\midrule
Convencional: no cubierto & 1 & 18 \\
Convencional: cubierto & 0 & 7 \\
\bottomrule
\end{tabular}
\end{table}

\subsection{Cobertura por bloque normativo}

La Figura~\ref{fig:bloques} desglosa la cobertura por los tres bloques normativos. En Buenas Pr\'acticas Agr\'icolas (BPA, Resoluci\'on 183 de AGROCALIDAD), la cobertura convencional ya alcanzaba el 71,4\,\% (5 de 7 criterios) y llega al 100\,\% tras el m\'etodo legal-first. En Bioseguridad (Resoluci\'on 0072), la cobertura convencional era del 33,3\,\% (2 de 6 criterios) y tambi\'en llega al 100\,\% tras el m\'etodo. El bloque con mayor vac\'io es la Ley Org\'anica de Protecci\'on de Datos Personales (LOPDP), donde la elicitaci\'on convencional no cubri\'o ning\'un criterio, 0\,\% (0 de 13), y el m\'etodo legal-first eleva la cobertura al 92,3\,\% (12 de 13).

\begin{figure}[htbp]
\centering
\includegraphics[width=0.8\textwidth]{../07_Datos/resultados/curva_o_barras.png}
\caption{Cobertura de criterios por bloque normativo, antes (convencional) y despu\'es (legal-first)}\label{fig:bloques}
\end{figure}

\subsection{El criterio sin cobertura}

El \'unico criterio que permanece sin cobertura tras aplicar el m\'etodo legal-first es C11 (bloque LOPDP, Art. 48), que exige la designaci\'on de un delegado de protecci\'on de datos cuando el tratamiento es a gran escala o involucra datos sensibles. El equipo document\'o en \texttt{Modelo\_Legal\_LOPDP.md} que ese criterio es condicional al volumen real de trabajadores y productores que gestione el sistema, un dato que debe fijarse junto con el administrador de la finca antes de evaluarlo como cubierto o no cubierto, y esa verificaci\'on qued\'o pendiente al cierre del an\'alisis.

\section{Discusi\'on}\label{sec:discusion}

Esta secci\'on interpreta los resultados en tres apartados, las implicaciones para la investigaci\'on, las implicaciones para la pr\'actica de la Ingenier\'ia de Requerimientos en dominios agroindustriales regulados, y los hallazgos contraintuitivos. La interpretaci\'on se mantiene dentro del alcance del estudio de caso \'unico y del marco legal ecuatoriano.

\paragraph{Implicaciones para la investigaci\'on.} Los resultados responden la pregunta de investigaci\'on de forma directa: la elicitaci\'on convencional, entrevistas y sesiones de validaci\'on con walkthrough, deja sin cubrir el 73,1\,\% de los criterios legales verificables del caso, y el enfoque legal-first de Amaral et al. \cite{amaral2021} cierra la mayor parte de ese vac\'io, con una diferencia significativa y sin ning\'un criterio que pierda cobertura. Esto extiende la validaci\'on emp\'irica de Amaral et al., hecha sobre el RGPD europeo y un contrato de tratamiento de datos, a un marco legal latinoamericano y a un dominio agroindustrial no cubierto por la literatura revisada en la Secci\'on~\ref{sec:trabajo-relacionado}, en l\'inea con lo se\~nalado por Torre et al. \cite{torre2021} y Negri-Ribalta et al. \cite{negriribalta2024} sobre la escasez de evidencia primaria fuera de Europa.

\paragraph{Implicaciones para la pr\'actica.} El hallazgo con mayor relevancia pr\'actica es el bloque de la LOPDP, donde la elicitaci\'on convencional no cubri\'o ning\'un criterio pero el m\'etodo legal-first alcanz\'o el 92,3\,\%. Esto sugiere que, en dominios donde el marco legal es reciente, como la LOPDP ecuatoriana vigente desde 2021, los usuarios entrevistados no articulan espont\'aneamente obligaciones legales que desconocen o que a\'un no forman parte de la pr\'actica habitual del sector, y que un m\'etodo que deriva requisitos directamente del texto legal es necesario, no solo complementario, para cerrar ese tipo de vac\'io.

\paragraph{Hallazgo contraintuitivo.} El bloque de BPA ya alcanzaba una cobertura convencional relativamente alta, 71,4\,\%, antes de aplicar el m\'etodo legal-first, lo que indica que las pr\'acticas agr\'icolas cotidianas que el personal ya sigue de forma expl\'icita se verbalizan con facilidad en las entrevistas, a diferencia de las obligaciones de protecci\'on de datos, que son m\'as abstractas y no est\'an incorporadas a la rutina de campo. Que ese mismo bloque, junto con Bioseguridad, cierre el 100\,\% de su cobertura con el m\'etodo legal-first, y que ning\'un criterio retroceda en ning\'un bloque, respalda la propiedad aditiva del m\'etodo, discutida por Amaral et al. \cite{amaral2021}: los requisitos derivados del texto legal se suman a los ya elicitados sin invalidarlos.

\section{Amenazas a la validez}\label{sec:amenazas}

\paragraph{Validez interna.} La asignaci\'on de cubierto o no cubierto de cada criterio legal es en parte subjetiva, ya que quien eval\'ua podr\'ia inclinarse a marcar m\'as criterios como cubiertos tras aplicar el m\'etodo legal-first. Para mitigarlo, la regla de decisi\'on se documenta por escrito antes de evaluar, un criterio se marca cubierto solo si existe al menos un requisito verificable y medible que lo satisface, y la asignaci\'on se contrasta con al menos un segundo revisor independiente. Adem\'as, el modelo de 26 criterios podr\'ia no capturar la totalidad de las obligaciones legales relevantes, por lo que cada criterio queda enlazado a su art\'iculo legal verificado y el modelo puede auditarse desde fuera.

\paragraph{Validez externa.} El estudio es un caso \'unico, por lo que los hallazgos pueden no generalizar a otras fincas o dominios. Esto se mitiga con documentaci\'on exhaustiva del contexto, siguiendo a Runeson y H\"ost \cite{runeson2009}. El modelo de cumplimiento cubre un marco legal espec\'ifico, y los resultados pueden no generalizar a otras jurisdicciones o cultivos. El alcance legal se declara de forma expl\'icita y se discute como limitaci\'on, con la misma l\'ogica de Amaral et al. \cite{amaral2021} aplicada originalmente al RGPD europeo.

\paragraph{Validez de constructo.} Los 26 criterios son un constructo definido por el equipo a partir de la lectura del texto legal y podr\'ian no capturar todos los aspectos de un cumplimiento verificable. Cada criterio se deriva de un art\'iculo legal expl\'icito y se documenta en la matriz de trazabilidad.

\paragraph{Validez de conclusi\'on.} El n\'umero de unidades comparadas, 26 criterios, es peque\~no y limita la potencia de la prueba de McNemar. Se reporta el intervalo de confianza del 95 por ciento de la diferencia de proporciones por bootstrap adem\'as del valor $p$, y se complementa con el an\'alisis descriptivo por bloque normativo y por fuente de evidencia.

\paragraph{Registro del protocolo.} El protocolo se registr\'o en OSF (https://osf.io/7cvhy) el 2 de agosto de 2026, con posterioridad al inicio de la recolecci\'on de datos el 20 de junio de 2026. En consecuencia, las hip\'otesis y el plan de an\'alisis no estaban preinscritos antes de la elicitaci\'on de campo, y los an\'alisis deben leerse como exploratorios, no confirmatorios. Esta desviaci\'on respecto de un registro previo ideal est\'a documentada en \texttt{06\_Experimento/osf\_deviations.tex} y no se disimula en ning\'un punto de este manuscrito.

\section{Conclusiones y trabajo futuro}\label{sec:conclusiones}

Este estudio de caso responde la pregunta de investigaci\'on con evidencia emp\'irica: en Agr\'icola Moreira, la elicitaci\'on convencional cubri\'o el 26,9\,\% de los 26 criterios legales de protecci\'on de datos y trazabilidad agroindustrial verificables, y el enfoque legal-first de Amaral et al. \cite{amaral2021} elev\'o esa cobertura al 96,2\,\%, una diferencia de 69,2 puntos porcentuales estad\'isticamente significativa ($\chi^2=16,06$, $p<0,001$) y sin retrocesos.

Las contribuciones de este trabajo son tres. Primero, un modelo conceptual verificable de cumplimiento legal para un sistema agroindustrial ecuatoriano, con 26 criterios trazables a su art\'iculo legal en tres bloques normativos. Segundo, evidencia emp\'irica de que el enfoque legal-first extiende su validaci\'on m\'as all\'a del RGPD europeo a un marco legal latinoamericano reciente y a un dominio agroindustrial de policultivo, subrepresentado en la literatura de cumplimiento legal en Ingenier\'ia de Requerimientos. Tercero, un caso donde el vac\'io de cobertura no es uniforme entre bloques normativos, sino que se concentra en el marco legal m\'as reciente y menos incorporado a la pr\'actica cotidiana del sector, lo que aporta un matiz emp\'irico sobre en qu\'e condiciones un m\'etodo legal-first resulta m\'as necesario.

Quedan al menos tres l\'ineas de trabajo futuro. Primera, verificar con el administrador de la finca el volumen real de trabajadores y productores gestionado por el sistema para cerrar la evaluaci\'on del \'unico criterio, C11, que permanece condicional. Segunda, replicar el modelo de 26 criterios en otras fincas de policultivo ecuatorianas para evaluar si el patr\'on de vac\'io concentrado en la LOPDP se sostiene fuera del caso \'unico aqu\'i estudiado. Tercera, extender el modelo conceptual a otros marcos legales latinoamericanos de protecci\'on de datos con estructura similar a la LOPDP, para evaluar la generalizaci\'on del m\'etodo legal-first m\'as all\'a de Ecuador.

\backmatter

\section*{Disponibilidad de datos y materiales}

El paquete de replicación, con las transcripciones anonimizadas, las respuestas del cuestionario sin columnas identificativas, el corpus de requisitos, la matriz de trazabilidad y los scripts de análisis, se deposita en Zenodo con DOI persistente 10.5281/zenodo.22307881 (https://doi.org/10.5281/zenodo.22307881), con licencia CC BY 4.0 para los datos y la documentación y licencia MIT para el código. El protocolo y sus desviaciones quedan registrados en OSF (https://osf.io/7cvhy). El repositorio de código fuente quedó archivado en Software Heritage con el identificador swh:1:snp:465aaeba1b5d8a07e1c7bca122fc8277812e825a. El material identificable, consentimientos originales, audios y videos, permanece cifrado y bajo custodia del docente responsable, y no forma parte del depósito abierto.

\section*{Uso de tecnolog\'ias asistidas por IA}

Se utilizo un modelo de lenguaje para apoyar la organizacion documental del proyecto y generar borradores de secciones no dependientes de datos empiricos, como la estructura del manuscrito y el trabajo relacionado a partir de referencias verificadas por el equipo, y para pulir redaccion ya escrita por el equipo. El modelo tambien redacto un borrador de las secciones de Resultados, Discusion y Conclusiones a partir unicamente de las cifras reales y ya calculadas por los scripts versionados de \texttt{07\_Datos/scripts/} sobre \texttt{07\_Datos/datos\_crudos/cobertura\_legal.csv}, sin generar, alterar ni inventar ningun dato, cifra o resultado. El equipo reviso y se aproprio de esa interpretacion antes del cierre de la entrega. El modelo no genero, en ningun momento, datos empiricos, cifras experimentales ni resultados estadisticos por si mismo.
\section*{Agradecimientos}

Los autores agradecen al personal de Agr\'icola Moreira que colabor\'o de forma voluntaria con las entrevistas y las sesiones de validaci\'on de este estudio.

\bibliography{referencias}

\end{document}
